% !TEX root = ../main.tex
% ==============================================================================
% PORTADA DE LA PRÁCTICA
% ==============================================================================
% Edita los datos del grupo, autores y fecha en la sección de configuración
% de main.tex (o directamente aquí si lo prefieres).
% ==============================================================================

\begin{titlepage}
    \thispagestyle{empty}

    % --- Elemento gráfico sutil (barra lateral institucional en TikZ) ---
    \begin{tikzpicture}[remember picture, overlay]
        % Banda vertical izquierda en rojo institucional URJC y gris oscuro
        \fill[urjcred] (current page.north west) rectangle ([xshift=0.65cm, yshift=-11.5cm]current page.north west);
        \fill[urjcdark] ([yshift=-11.5cm]current page.north west) rectangle ([xshift=0.65cm]current page.south west);
        % Línea fina decorativa paralela
        \draw[urjcred!35, line width=0.8pt] ([xshift=0.85cm]current page.north west) -- ([xshift=0.85cm]current page.south west);
    \end{tikzpicture}

    \vspace*{0.2cm}

    % --- Cabecera institucional ---
    \noindent
    \begin{minipage}[c]{0.18\textwidth}
        \IfFileExists{assets/logo_urjc.png}{%
            \includegraphics[width=\linewidth]{assets/logo_urjc.png}%
        }{%
            % Emblema vectorial minimalista inspirado en la URJC
            \begin{tikzpicture}[scale=0.55]
                % Corona (tres círculos + barra)
                \fill[urjcred] (0.25, 2.35) circle (0.16);
                \fill[urjcred] (0.85, 2.52) circle (0.22);
                \fill[urjcred] (1.45, 2.35) circle (0.16);
                \fill[urjcred] (0.12, 1.98) rectangle (1.58, 2.14);
                % Letra U estilizada
                \fill[urjcdark] (0.12, 1.80) -- (0.46, 1.80) -- (0.46, 0.72)
                    arc[start angle=180, end angle=360, radius=0.39]
                    -- (1.24, 1.80) -- (1.58, 1.80) -- (1.58, 0.72)
                    arc[start angle=360, end angle=180, radius=0.73] -- cycle;
            \end{tikzpicture}%
        }
    \end{minipage}%
    \hfill
    \begin{minipage}[c]{0.79\textwidth}
        {\Large\bfseries\color{urjcdark} Universidad Rey Juan Carlos}\\[0.25em]
        {\normalsize\scshape\color{urjcgray} Escuela Técnica Superior de Ingeniería Informática}\\[0.15em]
        {\normalsize\color{urjcred}\bfseries Grado en Matemáticas}
    \end{minipage}

    \vspace{0.8cm}
    \noindent{\color{urjcred}\rule{\linewidth}{1.2pt}}

    \vspace{2.6cm}

    % --- Etiqueta de asignatura y curso ---
    \noindent
    {\small\bfseries\color{urjcred}\MakeUppercase{\asignatura} \quad $\cdot$ \quad Curso \curso}

    \vspace{0.6cm}

    % --- Título principal ---
    \noindent
    {\fontsize{30}{36}\selectfont\bfseries\color{urjcdark} \titulopractica \par}

    \vspace{0.5cm}

    % --- Subtítulo descriptivo ---
    \noindent
    {\large\color{urjcgray} \subtitulopractica \par}

    \vspace{1.2cm}
    \noindent{\color{urjcdark!25}\rule{0.35\linewidth}{1.5pt}}

    \vfill

    % --- Bloque de metadatos (Autores, Grupo, Fecha) ---
    \noindent
    \begin{tikzpicture}
        \node[
            fill=urjcbglight,
            draw=urjcdark!15,
            rounded corners=4pt,
            inner sep=16pt,
            text width=\dimexpr\linewidth-32.8pt\relax,
            anchor=north west
        ] {
            \renewcommand{\arraystretch}{1.45}
            \begin{tabular}{@{}l@{\hspace{1.2em}}p{0.62\linewidth}@{}}
                \textbf{\color{urjcdark} Integrantes del grupo:} &
                \begin{minipage}[t]{\linewidth}\raggedright \autores \end{minipage} \\[0.4em]
                \textbf{\color{urjcdark} Grupo:}                 & \grupoclase \\
                \textbf{\color{urjcdark} Fecha de entrega:}      & \fechaentrega \\
            \end{tabular}
        };
    \end{tikzpicture}

    \vspace{0.8cm}

    % --- Pie de portada ---
    \begin{center}
        {\footnotesize\color{urjcgray} Memoria de Práctica elaborada en \LaTeX\ --- Móstoles, Madrid}
    \end{center}

\end{titlepage}
